
% Valerie Fong — one-page quantitative research / data science resume
% Build: pdflatex resume.tex
\documentclass[10pt,letterpaper]{article}

\usepackage[T1]{fontenc}
\usepackage[utf8]{inputenc}
\usepackage{lmodern}
\usepackage[top=0.45in,bottom=0.45in,left=0.55in,right=0.55in]{geometry}
\usepackage{xcolor}
\usepackage{enumitem}
\usepackage{titlesec}
\usepackage[hidelinks,pdfusetitle]{hyperref}
\usepackage{microtype}
\hyphenpenalty=10000
\exhyphenpenalty=10000

% Improve PDF text extraction and ATS parsing
\input{glyphtounicode}
\pdfgentounicode=1

\hypersetup{
  pdftitle={Valerie Fong Resume},
  pdfauthor={Valerie Fong}
}

\definecolor{ink}{HTML}{000000}
\definecolor{muted}{HTML}{000000}
\definecolor{accent}{HTML}{000000}
\definecolor{rule}{HTML}{000000}
\definecolor{sep}{HTML}{000000}

\pagestyle{empty}
\setlength{\parindent}{0pt}
\setlength{\parskip}{0pt}
\renewcommand{\baselinestretch}{1}
\AtBeginDocument{\fontsize{10.2}{13}\selectfont\color{ink}}

% Section headings
\titleformat{\section}
  {\fontsize{11.5}{13}\selectfont\bfseries\color{accent}}
  {}{0pt}{\MakeUppercase}
  [{\vspace{1pt}\color{rule}\titlerule[0.5pt]}]
\titlespacing*{\section}{0pt}{7pt}{3.5pt}

\setlist[itemize]{
  leftmargin=12pt,
  labelsep=5pt,
  topsep=1pt,
  itemsep=1.2pt,
  parsep=0pt,
  partopsep=0pt,
  label=\textbullet
}

% \entry{left}{date}
\newcommand{\entry}[2]{%
  \noindent#1\hfill
  {\fontsize{9.8}{12}\selectfont\color{muted}#2}\par
}
\newcommand{\sub}[1]{{\itshape\color{muted}#1\par}}
\newcommand{\gap}{\vspace{4pt}}
\newcommand{\cbar}{{\color{sep}\enspace|\enspace}}

\begin{document}

\begin{center}
{\fontsize{21}{24}\selectfont Valerie Fong}\par\vspace{2pt}
{\fontsize{9.8}{12}\selectfont\color{muted}%
\href{mailto:valfong@ucdavis.edu}{valfong@ucdavis.edu}\cbar
505-850-3795\cbar
\href{https://valeriefong.github.io}{valeriefong.github.io}\cbar
\href{https://github.com/valeriefong}{github.com/valeriefong}\cbar
\href{https://linkedin.com/in/valfong}{linkedin.com/in/valfong}}
\end{center}
\vspace{-8pt}

\section*{Education}

\entry{\textbf{University of California, Davis}, Ph.D. in Applied Mathematics}
{Expected Dec 2027}
\sub{Dissertation: Temporal Computation through Dynamically Critical Networks. Advisor: Rishidev Chaudhuri.}

\gap
\entry{\textbf{University of California, Davis}, M.S. in Applied Mathematics}
{2024}

\gap
\entry{\textbf{University of New Mexico}, B.S. in Statistics and Economics, minor in Mathematics}
{2022}
\sub{Phi Beta Kappa}
\gap
\textbf{Fellowships:} NSF Graduate Research Fellowship; National GEM Consortium Ph.D. Fellowship; UC Davis Hazel B. Jacoby Fellowship

\section*{Research Experience}

\entry{\textbf{University of California, Davis}, Graduate Student Researcher}
{2022 -- Present}
\begin{itemize}
  \item Developed a Python/PyTorch simulation pipeline with SLURM job arrays to study Lyapunov spectra in 3 recurrent network models, scaling to $N=2{,}000$ across 1,060 parameter combinations and approximately 25,000 simulation runs.
  \item Adapted and extended a published Lyapunov-spectrum solver using the QR/Benettin method, fixed-step RK4 integration, and tangent dynamics; validated the pipeline against the reference implementation and assessed numerical convergence through paired comparisons of step size, orthonormalization interval, simulation duration, precision, and network size.
  \item Executed the resulting computational study on an HPC cluster, using approximately 70,000 CPU core-hours to characterize how network size, model architecture, and connectivity gain shape Lyapunov spectra.
  \item Explored Dynamic Mode Decomposition on high-dimensional neural recordings, working with the experimental neuroscientists who collected the data to understand its biological and experimental context.
\end{itemize}

\gap
\entry{\textbf{University of New Mexico}, Ronald E. McNair Scholar and ASSURE Fellow}
{2021 -- 2022}
\begin{itemize}
  \item Applied variance-based Sobol' global sensitivity analysis to quantify how uncertainty in eight model parameters influences Hodgkin--Huxley neuron model outputs across hundreds of thousands of simulations.
  \item Presented research findings at three national McNair research conferences.
\end{itemize}

\section*{Professional Experience}

\entry{\textbf{iRhythm Technologies}, Data Science Intern}
{Jun -- Aug 2022}
\begin{itemize}
  \item Built and benchmarked time-series demand forecasting models, including ARIMA, convolutional neural networks, empirical mode decomposition, and seasonal-trend decomposition.
  \item Compared forecasting approaches across multiple error metrics and applied k-means clustering and principal component analysis (PCA) to investigate patterns in demand data.
\end{itemize}

\gap
\entry{\textbf{Sandia National Laboratories}, Year-Round R\&D Intern}
{2019 -- 2021}
\begin{itemize}
  \item Developed and maintained a dashboard of safety, security, quality, and environmental metrics to support leadership decision-making.
  \item Earned Lean Six Sigma Green Belt certification; held a U.S. Department of Energy Q clearance.
\end{itemize}

\section*{Teaching}

\entry{\textbf{Private Tutor}, Calculus and Ordinary Differential Equations}
{2026}

\gap
\entry{\textbf{Peer Learning Facilitator}, Advanced Data Analysis I \& II (University of New Mexico)}
{2021 -- 2022}

\section*{Technical Skills}\raggedright

\textbf{Programming:} Python (NumPy, SciPy, pandas, scikit-learn, Matplotlib, PyTorch), R, MATLAB, Git, LaTeX, HPC/SLURM\par\vspace{2pt}

\textbf{Statistics and Modeling:} Time-series forecasting, statistical modeling, global sensitivity analysis, clustering, PCA\par\vspace{2pt}

\textbf{Numerical Methods:} Numerical ODE integration, RK4, eigenspectral methods, Lyapunov exponents, QR/Benettin method, Dynamic Mode Decomposition, large-scale simulation

\end{document}
